\usepackage{lmodern} % font
\usepackage[english]{babel}
\usepackage{csquotes}
\usepackage{lipsum, kantlipsum}
\usepackage{etoolbox}

% For Title styling
\usepackage{titlesec}
% MD: \raggedright sui titoli. Senza, titlesec li giustifica e quelli lunghi
% sforano il margine destro nel layout finale A4: "Electromagnetic Trapping of
% Neutral Atoms" di 39 pt, "From Conventional to Hollow-Core Optical Fibers"
% di 20 pt, "Coupling to the Motional ..." di 14 pt. Cosi' vanno tutti a zero.
% NB: l'opzione di pacchetto \usepackage[raggedright]{titlesec} NON basta,
% perche' \titleformat* la sovrascrive: deve stare dentro l'argomento.
\titleformat*{\section}{\LARGE\bfseries\raggedright}
\titleformat*{\subsection}{\Large\bfseries\raggedright}
\titleformat*{\subsubsection}{\large\bfseries\raggedright}

% Margins and layout
\usepackage[margin=2.5cm]{geometry}
\usepackage{setspace}
\onehalfspacing % Increase the line spacing
\setlength{\headheight}{14.5pt}

% Graphics and figures
\usepackage{graphicx, adjustbox}
\usepackage{float}
\usepackage{subcaption}
\graphicspath{{figures/}}

\usepackage{wrapfig}

% Mathematics
\usepackage{amsmath,amssymb,amsthm}
% Physics
\usepackage{braket}

% For organizing the work in subfiles
\usepackage{subfiles}

% Bibliography
\usepackage[backend=biber,style=numeric,sorting=none]{biblatex}
\addbibresource{\subfix{bibliography.bib}}
\DefineBibliographyStrings{english}{%
  andothers = {\mkbibemph{et~al.\adddot}},
}

% Hyperlinks
\usepackage{xr-hyper}
\usepackage[colorlinks=true,linkcolor=black,citecolor=black,urlcolor=black]{hyperref}

% Headers and footers
\usepackage{fancyhdr}
\usepackage{extramarks}
\pagestyle{fancy}
\fancyhf{}
\fancyhead[LE,RO]{\thepage} % Left Even, Right Odd number page
\fancyhead[RE]{\leftmark}   % Right Even, left mark (Chapter title)
\fancyhead[LO]{\rightmark}  % Left Odd, right mark (Section name)


% No indentation on new paragraphs
\setlength{\parindent}{0pt}
% MD: con \parindent a zero e nessun \parskip i capoversi erano
% tipograficamente indistinguibili, si leggevano come un blocco unico.
% Servendo l'uno o l'altro, qui si aggiunge lo spazio verticale.
\setlength{\parskip}{0.5\baselineskip plus 2pt minus 1pt}

% For numbering also subsubsections
\setcounter{secnumdepth}{3}
\setcounter{tocdepth}{3}

% Tikz 
\usepackage{tikz}
\usetikzlibrary{calc, positioning, arrows.meta}
\usepackage{xcolor}
\usepackage{optikz}

% Image caption in itlalicum
\usepackage[textfont=it]{caption}

% Itemize stuff
\usepackage{enumitem}
% Better Tabels
\usepackage{booktabs, multirow, array, makecell}
\newcolumntype{C}[1]{>{\centering\arraybackslash}p{#1}}

% review by MICHELANGELO
%\geometry{paperwidth=27cm, asymmetric, left=2.5cm, right=8.5cm, marginparwidth=7cm, marginparsep=5mm}
%\usepackage[colorinlistoftodos, textsize=footnotesize]{todonotes}

% scientific notes
%\providecommand{\MDn}[1]{\todo[color=brown!20]{\textbf{MD notes:} #1}}

% questions
%\providecommand{\MDq}[1]{\todo[inline, color=blue!15]{\textbf{MD question:} #1}}

% suggestions
%\providecommand{\MDs}[1]{\todo[noline, color=red!25]{\textbf{MD suggestion:} #1}}

% typographical issue
%\providecommand{\MDt}[1]{\todo[color=orange!30]{\textbf{MD typography:} #1}}

% Underlining text
%\usepackage{ulem}

% Geometry setup for wide margins
\geometry{paperwidth=30cm, asymmetric, left=1.5cm, right=11.5cm, marginparwidth=10cm, marginparsep=5mm}

% Required packages
\usepackage{soul} % for multiline highlighting (\hl)
\usepackage[normalem]{ulem} % [normalem] prevents ulem from changing \emph to underline
\usepackage[colorinlistoftodos, textsize=footnotesize]{todonotes}

% --- Custom Color Setup for soul Highlights ---
% \leavevmode: senza, \hl fallisce con "Reconstruction failed" quando una nota
% apre il capoverso subito dopo un \end{equation} (soul parte in modo verticale).
\newcommand{\hlcolor}[2]{\leavevmode{\sethlcolor{#1}\hl{#2}}}

% --- 2-Argument Review Commands: \Command{Text to highlight}{Comment} ---

% 1. Scientific notes (brown highlight + margin note with connecting line)
\newcommand{\MDn}[2]{%
  \hlcolor{brown!25}{#1}%
  \todo[color=brown!20]{\textbf{MD note:} #2}%
}

% 2. Questions (blue highlight + margin note)
\newcommand{\MDq}[2]{%
  \hlcolor{blue!15}{#1}%
  \todo[color=blue!15]{\textbf{MD question:} #2}%
}

% 3. Phrasing / Suggestions (red/pink highlight + margin note)
\newcommand{\MDs}[2]{%
  \hlcolor{red!20}{#1}%
  \todo[color=red!25]{\textbf{MD suggestion:} #2}%
}

% 4. Typographical / Grammar (wavy underline from ulem + margin note)
\newcommand{\MDt}[2]{%
  \uwave{#1}%
  \todo[color=orange!30]{\textbf{MD typo:} #2}%
}

% ATTENZIONE: non usare \MDn / \MDq / \MDs / \MDt / \MDdel dentro l'argomento di
% \chapter, \section, \subsection ... . Sono argomenti mobili: la nota finisce
% verbatim nell'indice, hyperref segnala "Token not allowed in a PDF string" e
% il \todo viene eseguito due volte (una dall'indice, una dal titolo).
% Mettere la nota subito DOPO il titolo, nella forma \MDt{}{...}.
% Per lo stesso motivo, non usare \ref / \cite dentro \MDn ... \MDdel: sono
% comandi fragili, e sia \hl (soul) sia \todo (che scrive nel file .tdo)
% li rompono. Citare l'etichetta come testo, es. \texttt{eq:dipole\_force}.

% 1b. Variante di \MDn senza soul. \hl fallisce con "Reconstruction failed" in
% certe posizioni (paragrafi accorciati da wrapfig, o a cavallo di
% un'interruzione di pagina): stessa semantica di \MDn, ma il testo viene
% sottolineato invece che evidenziato. Usare solo dove \MDn non compila.
\newcommand{\MDnx}[2]{%
  \uline{#1}%
  \todo[color=brown!20]{\textbf{MD note:} #2}%
}

% 5. Direct Deletion (strikethrough from ulem + margin note)
\newcommand{\MDdel}[2]{%
  \sout{#1}%
  \todo[color=gray!20]{\textbf{MD delete:} #2}%
}